\PassOptionsToPackage{dvipsnames,svgnames,table}{xcolor}
\documentclass[10pt]{article}
\usepackage[a4paper,margin=22mm]{geometry}
\usepackage[no-math]{fontspec}\setmainfont{Latin Modern Roman}
\usepackage{amsmath,amssymb,amsthm,mathtools,graphicx,xcolor,hyperref,xurl,xspace,ifthen,enumerate}
\usepackage[shortlabels]{enumitem}
\setlength{\emergencystretch}{4em}
\newtheorem{theorem}{Theorem}[section]
\newtheorem{theo}[theorem]{Theorem}
\newtheorem{thm}[theorem]{Theorem}
\newtheorem*{thm*}{Theorem}
\newtheorem{lemma}[theorem]{Lemma}
\newtheorem{lemm}[theorem]{Lemma}
\newtheorem{prop}[theorem]{Proposition}
\newtheorem{coro}[theorem]{Corollary}
\theoremstyle{definition}
\newtheorem{defi}[theorem]{Definition}
\newtheorem{rema}[theorem]{Remark}
\newtheorem{exam}[theorem]{Example}
\providecommand{\eg}{e.g.\ }
\providecommand{\ie}{i.e.\ }
\providecommand{\resp}{resp.\ }
\providecommand{\cf}{cf.\ }
\providecommand{\longthanks}[1]{\par\smallskip\noindent #1\par\smallskip}
\usepackage{tikz}

\newcommand{\edgedir}{\mathbin{\tikz [semithick, baseline=-0.2ex,-latex, ->] \draw [->] (0pt,0.4ex) -- (1em,0.4ex);}} %
\newcommand{\edgedirback}{\mathbin{\tikz [semithick, baseline=-0.2ex,-latex, ->] \draw [<-] (0pt,0.4ex) -- (1em,0.4ex);}} %
\newcommand*{\mk}{\mkern -1mu}
\newcommand*{\Mk}{\mkern -2mu}
\newcommand*{\mK}{\mkern 1mu}
\newcommand*{\MK}{\mkern 2mu}

\hypersetup{urlcolor=purple, linkcolor=blue, citecolor=red}


\newcommand*{\romanenumi}{\renewcommand*{\theenumi}{\roman{enumi}}}
\newcommand*{\Romanenumi}{\renewcommand*{\theenumi}{\Roman{enumi}}}
\newcommand*{\alphenumi}{\renewcommand*{\theenumi}{\alph{enumi}}}
\newcommand*{\Alphenumi}{\renewcommand*{\theenumi}{\Alph{enumi}}}
\begin{document}
\section*{1508. Hurwitz normalization to paired excess reflections at the simple-length threshold in finite-rank Coxeter groups}
Wegener, Patrick; Yahiatene, Sophiane. \emph{A note on non-reduced reflection factorizations of Coxeter elements}. Algebraic Combinatorics 3 (2020), publisher version, DOI 10.5802/alco.99. \url{https://alco.centre-mersenne.org/articles/10.5802/alco.99/}. CC BY 4.0, \url{https://creativecommons.org/licenses/by/4.0/}.
\paragraph{Selected problem and scope (editorial).} Exact original Lemma 2.3 for any finite-rank Coxeter system, arbitrary w with simple-generator length m and any reflection factorization of length m+2k. Complete original rank-two peak replacement Proposition 2.2, its proof, Hurwitz action, reflection subgroup conventions and the full descent/induction are included. The prefix is not asserted to be reflection-reduced. Original Theorem 1.2 remains unselected contextual motivation; neither its final proof nor orbit-classification conclusion is counted. Named Dyer subgroup/coset Bruhat results remain established prerequisites; no new proof is generated.
\paragraph{Difficulty (editorial).} After the stated Dyer reflection-subgroup and coset Bruhat-graph prerequisites, the selected normalization still requires an explicit rank-two peak repair that lowers the global simple-length potential and a terminating Hurwitz descent for arbitrary finite-rank Coxeter groups. The author then uses the identity-starting peak-free path obstruction to force a repeated reflection and completes the paired-suffix induction, at native179–245. This uniform non-reduced factorization construction exceeds ordinary qualification exercises in Coxeter relations; conservative H2 is proposed.
\paragraph{Formatting (editorial).} Original mathematical text, formulas, diagrams and proof order are retained. The publisher class is replaced by a standalone article scaffold, with original numbering restored by wrapper counters where needed. Non-rendered original comment prose is excluded while every percent newline guard is retained. Bibliography facts are reproduced under their original citation keys from the exact archived publisher HTML. No publisher class or upstream script is executed, and no proof or formula is generated.
\clearpage

\section{Introduction}\label{sec:intro}


Given a Coxeter system $(W,S)$ with set of reflections $T$, the braid group (\eg see~\cite{BDSW14} for a definition) acts on \emph{reflection factorizations} of a given element $w \in W$, that is it acts on tuples $(t_1, \ldots , t_m) \in T^m$ of reflections such that $w=t_1 \cdots t_m$. This action is called \emph{Hurwitz action}. A standard braid group generator $\sigma_i$ (resp. its inverse $\sigma_i^{-1}$) acts by a \emph{Hurwitz move} on a reflection factorization:
\begin{align*}
\sigma_i (t_1, \ldots , t_{i-1}, t_i, t_{i+1}, t_{i+2}, \ldots , t_n) & = (t_1, \ldots , t_{i-1}, t_{i+1}^{t_i}, t_{i}, t_{i+2}, \ldots , t_n),\\
\sigma_i^{-1} (t_1, \ldots , t_{i-1}, t_i, t_{i+1}, t_{i+2}, \ldots , t_n) & = (t_1, \ldots , t_{i-1}, t_{i+1}, t_{i}^{t_{i+1}}, t_{i+2}, \ldots , t_n),
\end{align*}
where we use the notation $g^h:=hgh^{-1}$ for conjugation. 

We call a reflection factorization $(t_1, \ldots , t_m)$ of an element $w \in W$ \emph{reduced} if $w$ cannot be written as a product of less than $m$ reflections. It has been first observed by Deligne~\cite{Del} that this action is transitive on reduced reflection factorizations of a \emph{Coxeter element} if $W$ is finite. The first published proof is due to Bessis~\cite[Proposition~1.6.1]{Bes03}. Igusa and Schiffler showed that this statement is true for every Coxeter group~\cite[Theorem~1.4]{IS10}. 

The question of how these results extend to non-reduced reflection factorizations has been first addressed by Lewis and Reiner.

\begin{thm*}
[{Lewis--Reiner,~\cite[Theorem~1.1]{LR16}}] In a finite real reflection group, two reflection factorizations of a Coxeter element lie in the same Hurwitz orbit if and only if they share the same multiset of conjugacy classes.
\end{thm*}

Their proof makes heavy use of the following remarkable result for finite Coxeter groups.
\begin{lemma}[{Lewis--Reiner,~\cite[Corollary~1.4]{LR16}}] \label{lem:LewisR}
Let $(W,S)$ be a finite Coxeter system and $w\in W$. Then every reflection factorization of $w$ lies in the same Hurwitz orbit of some reflection factorization $(t_1, \ldots , t_m)$ of $w$ such that $(t_1, \ldots , t_{\ell})$ is a reduced reflection factorization of $w$ for some $\ell \leq m$ and 
\[
t_{\ell +1} = t_{\ell+2}, ~t_{\ell+3} = t_{\ell+4}, \ldots, ~t_{m-1}=t_m.
\]
\end{lemma}

This result is proved by a case-by-case analysis and seems not to extend to infinite Coxeter groups in general. We give a case-free proof of a similar (but weaker) result for all Coxeter groups of finite rank (see Lemma~\ref{lem:same_refl}). In this way, we obtain that the result of Lewis and Reiner extends to all Coxeter groups of finite rank which provides a positive answer to a question of Lewis and Reiner~\cite[Question~6.2]{LR16}.

\begin{thm} \label{thm:Main}
Let $(W,S)$ be a Coxeter system of finite rank. Then two reflection factorizations of a Coxeter element in $W$ lie in the same Hurwitz orbit if and only if they share the same multiset of conjugacy classes.
\end{thm}



\section{The proof}
Throughout this note let $(W,S)$ be a \emph{Coxeter system} of finite rank $n \in \mathbb{N}$ with set of \emph{reflections} $T=\{ wsw^{-1} \mid w \in W,~s \in S\}$. All necessary definitions and facts about Coxeter groups we will use are covered by~\cite{BDSW14, Dye90, Dye91, Hu90}.

A subgroup $W'$ of $W$ is called a \emph{reflection subgroup} if $W'= \langle W' \cap T \rangle$. Each reflection subgroup $W'$ admits a canonical set of generators $\chi(W')$ such that $(W', \chi(W'))$ is a Coxeter system and the set of reflections for $(W', \chi(W'))$ is given by $W' \cap T = \bigcup_{w \in W'} w \chi(W') w^{-1}$ (see~\cite[(3.3)~Theorem]{Dye90}). A reflection subgroup of the form $\langle I \rangle$ for some $I \subseteq S$, is called \emph{parabolic subgroup}.

Let $S =\{ s_1, \ldots , s_n \}$. For each permutation $\pi$ of the numbers $\{1, \ldots, n\}$, the element $c= s_{\pi(1)} \cdots s_{\pi(n)}$ is called a \emph{Coxeter element}. A Coxeter element of a parabolic subgroup is called \emph{parabolic Coxeter element}.

We denote by $\ell_S$ (\resp $\ell_T$) the length function on $W$ with respect to the generating set $S$ (\resp $T$).



\begin{defi} \label{def:BruhatGraph}
We define the \emph{Bruhat graph} of $(W,S)$ to be the directed graph $\Omega_{(W,S)}$ on vertex set $W$ with a directed edge from $x$ to $y$ if there exists $t \in T$ such that $y=xt$ and $\ell_S(x) < \ell_S(y)$.

Moreover, we denote by $\overline{\Omega}_{(W,S)}$ the corresponding undirected graph and for a subset $X \subseteq W$ we denote by $\Omega_{(W,S)}(X)$ the full subgraph of $\Omega_{(W,S)}$ on the vertex set $X$.

Note that $\overline{\Omega}_{(W,S)}$ is exactly the Cayley graph of $W$ with generating set $T$. 
\end{defi}

We use the notation $(t_1, \ldots, t_m) \sim (r_1, \ldots , r_m)$ to indicate that both tuples lie in the same orbit under the Hurwitz action.

\medskip
The following fact is already part of the proof of~\cite[Proposition~2.2]{BDSW14}. For sake of completeness we include a proof (which can also be found in the first author's Ph.D. thesis~\cite[Proposition~2.3.6]{Weg17}).

\begin{prop} \label{prop:BruhatGraphDihedral}
Let $w \in W$ and $t_1,t_2 \in T$ with $t_1 \neq t_2$ such that 
\[ 
w \edgedir wt_1 \edgedirback wt_1t_2
\]
in $\Omega_{(W,S)}$. Then there exist $t_1', t_2 ' \in \langle t_1, t_2 \rangle \cap T$ with $(t_1, t_2) \sim (t_1', t_2')$ such that one of the following cases hold:
\begin{enumerate}
\item\label{prop2.2_1} $ w \edgedir wt_1' \edgedir wt_1't_2'= wt_1t_2$;
\item\label{prop2.2_2} $ w \edgedirback wt_1' \edgedirback wt_1't_2'= wt_1t_2$;
\item\label{prop2.2_3} $ w \edgedirback wt_1' \edgedir wt_1't_2'= wt_1t_2$;
\end{enumerate}
Furthermore, in each of these cases we have $\ell_S(wt_1')<\ell_S(wt_1)$.
\end{prop}

\begin{proof}
Consider the dihedral reflection subgroup $W':= \langle t_1, t_2 \rangle$ and let $S':= \chi(W')$. We have $w, wt_1, wt_1t_2 \in wW'$. Therefore we just have to consider the coset $wW'$ to prove the claim. By~\cite[(1.4)~Theorem]{Dye91} we have 
\[
\Omega_{(W,S)}(W') \cong \Omega_{(W,S)}(wW') \cong \Omega_{(W',S')},
\]
where $(W',S')$ is dihedral and we can check the claim there directly. Any reflection (element of odd $S'$-length) of $W'$ is joined by an edge to a rotation (element of even $S'$-length) which in $\Omega_{(W',S')}$ is oriented towards the element of greater $S'$-length. For $x \in W'$ there are three possible situations:
\begin{itemize}
\item $\ell_{S'}(x)< \ell_{S'}(xt_1t_2)$
\item $\ell_{S'}(x)> \ell_{S'}(xt_1t_2)$
\item $\ell_{S'}(x)= \ell_{S'}(xt_1t_2)$ (in particular $x \neq e$ since $t_1 \neq t_2$).
\end{itemize}
We can choose $t_1', t_2' \in W' \cap T$ with $t_1't_2' = t_1t_2$ in the three situations such that we have one of the following situations:
\begin{itemize}
\item $x \edgedir xt_1' \edgedir xt_1' t_2'$
\item $x \edgedirback xt_1' \edgedirback xt_1' t_2'$
\item $x \edgedirback xt_1' \edgedir xt_1' t_2'$
\end{itemize}
To see this, note that $x$ and $xt_1t_2$ are both either reflections or rotations. Therefore both are either of odd or even $S'$-length. Thus $\ell_{S'}(x) < \ell_{S'}(xt_1t_2)$ implies $\ell_{S'}(x) +2 \leq \ell_{S'}(xt_1t_2)$ and we find $t_1'$ with $\ell_{S'}(x) < \ell_{S'}(xt_1') < \ell_{S'}(xt_1t_2)$. By setting $t_2' := t_1't_1t_2$ we obtain $x \edgedir xt_1' \edgedir xt_1' t_2'$ and $t_1't_2'=t_1t_2$. The remaining cases are similar.

It is easy to see that the Hurwitz orbit of $(t_1,t_2)$ is the set of all pairs $(r_1,r_2)$ of reflections of $W'$ such that $t_1t_2=r_1r_2$ (see also~\cite{BDSW14}). Hence we have $(t_1, t_2) \sim (t_1', t_2')$.

It remains to show that $\ell_S(wt_1')<\ell_S(wt_1)$ in each of the cases~\eqref{prop2.2_1}-\eqref{prop2.2_3}. Since the initial path is of the form $w \edgedir wt_1 \edgedirback wt_1t_2$, we have:
\begin{enumerate}[(\roman*)]
\item\label{proof_prop2.2_i} %
$\ell_S(w) < \ell_S(wt_1)$, and
\item\label{proof_prop2.2_ii} %
$\ell_S(wt_1t_2) < \ell_S(wt_1) $.
\end{enumerate}
In case~\eqref{prop2.2_1} we have an edge $wt_1' \edgedir wt_1't_2'= wt_1t_2$, thus 
\[
\ell_S(wt_1') < \ell_S(wt_1't_2') = \ell_S(wt_1t_2) \stackrel{\ref{proof_prop2.2_ii}}{<} \ell_S(wt_1).
\]
In cases~\eqref{prop2.2_2} and~\eqref{prop2.2_3} we have an edge $w \edgedirback wt_1'$, thus 
\[
\ell_S(wt_1') < \ell_S(w) \stackrel{\ref{proof_prop2.2_i}}{<} \ell_S(wt_1).\qedhere
\]
\end{proof}




\begin{lemma}\label{lem:same_refl}
Let $w\in W$ with $\ell_{S}(w)=m$ and $w=t_{1}\cdots t_{m+2k}$ with $t_{i}\in T$ for $1\leq i \leq m+2k$ and some $k \in \mathbb{Z}_{\geq 0}$. Then there exists a braid $\sigma \in \mathcal{B}_{m+2k}$ such that 
\[
\sigma(t_{1},\ldots,t_{m+2k})=(r_{1},\ldots,r_{m},r_{i_1},r_{i_1}, \ldots , r_{i_k},r_{i_k}).
\]
\end{lemma}

\begin{proof}
We proceed by induction on $k$. The case $k=0$ is trivially satisfied. Therefore let $k\geq 1$ and assume that all factorizations in $\mathcal{B}_{m+2k}(t_{1},\ldots,t_{m+2k})$ consist of pairwise different factors. Consider the path of $\Omega_{(W,S)}$ starting in $e$ and ending in $w$ induced by $(t_{1},\ldots,t_{m+2k})$. Then Proposition~\ref{prop:BruhatGraphDihedral} allows us to replace successively the parts of the path of shape $\star \edgedir \star \edgedirback \star$ by
\[
\star \edgedir \star \edgedir \star, ~\star \edgedirback \star \edgedirback \star, ~\text{or } \star \edgedirback \star \edgedir {} \star
\] 
only using the Hurwitz action. The latter is possible since the reflections of the factorizations in the Hurwitz orbit are pairwise different. Since by Proposition~\ref{prop:BruhatGraphDihedral} each replacement reduces the sum of the length of the vertices, eventually we get after finitely many replacements a path of the form
\[
e\edgedirback t'_{1}\edgedirback t'_{1}t'_{2}\edgedirback \ldots \edgedirback t'_{1}t'_{2}\cdots t'_{p}\edgedir t'_{1}t'_{2}\cdots t'_{p}t'_{p+1}\edgedir \ldots \edgedir t'_{1}\cdots t'_{m+2k}=w
\]
with $t'_{i}\in T$ for $1\leq i \leq m+2k$, that is, the path is first decreasing, then increasing. Since the path starts with $e$, it holds $p=0$ and therefore it has no decreasing part. Altogether the initial path can be transformed to
\[
e \edgedir t'_{1}\edgedir t'_{1}t'_{2}\edgedir \ldots \edgedir t'_{1}\cdots t'_{m+2k}=w
\]
by using the Hurwitz action. Since the length of $w$ is $m$ and $k\geq 1$, there cannot be an increasing path of length $m+2k$, so we arrive at a contradiction. Thus there exists a factorization $(t_{1}',\ldots,t_{m+2(k-1)}',r_{i_{k}},r_{i_{k}})$ in $\mathcal{B}_{m+2k}(t_{1},\ldots,t_{m+2k})$. From the induction hypothesis follows 
\[
(t_{1}',\ldots,t_{m+2(k-1)}')\sim (r_{1},\ldots,r_{m},r_{i_{1}},r_{i_{1}},\ldots,r_{i_{k-1}},r_{i_{k-1}})
\] 
and the latter yields the assumption.
\end{proof}



\begin{rema}
Note that $\ell_T(w) \leq \ell_S(w)$ for all $w \in W$. Therefore the reflection factorization $w = r_1 \cdots r_m$ obtained by Lemma~\ref{lem:same_refl} does not have to be a reduced reflection factorization. This is the main difference compared with the key argument Lemma~\ref{lem:LewisR} in the proof of Lewis and Reiner. However, by~\cite[Lemma~2.1]{BDSW14} we have $\ell_S(w)=\ell_T(w)$ for an element $w \in W$ if and only if $w$ is a parabolic Coxeter element. Therefore, if $w$ is a parabolic Coxeter element, then Lemma~\ref{lem:same_refl} generalizes Lemma~\ref{lem:LewisR}. In particular, a reflection factorization of a parabolic Coxeter element can be reduced by just using Hurwitz moves and deleting matching neighbors.
\end{rema}
\begin{thebibliography}{99}
\bibitem[1]{BDSW14} Baumeister, Barbara; Dyer, Matthew; Stump, Christian; Wegener, Patrick A note on the transitive Hurwitz action on decompositions of parabolic Coxeter elements, Proc. Am. Math. Soc., Ser. B, Volume 1 (2014), pp. 149-154.
\bibitem[2]{Bes03} Bessis, David The dual braid monoid, Ann. Sci. Éc. Norm. Supér. (4), Volume 36 (2003) no. 5, pp. 647-683.
\bibitem[3]{Del} Deligne, Pierre Letter to E. Looijenga (1974) \url{http://homepage.rub.de/christian.stump/Deligne_Looijenga_Letter_09-03-1974.pdf}.
\bibitem[4]{Dye90} Dyer, Matthew Reflection subgroups of Coxeter systems, J. Algebra, Volume 135 (1990) no. 1, pp. 57-73.
\bibitem[5]{Dye91} Dyer, Matthew On the “Bruhat graph” of a Coxeter system, Compos. Math., Volume 78 (1991) no. 2, pp. 185-191.
\bibitem[6]{Hu90} Humphreys, James E. Reflection groups and Coxeter groups, Cambridge Studies in Advanced Mathematics, 29, Cambridge University Press, Cambridge, 1990, xii+204 pages.
\bibitem[7]{IS10} Igusa, Kiyoshi; Schiffler, Ralf Exceptional sequences and clusters, J. Algebra, Volume 323 (2010) no. 8, pp. 2183-2202.
\bibitem[8]{LR16} Lewis, Joel Brewster; Reiner, Victor Circuits and Hurwitz action in finite root systems, New York J. Math., Volume 22 (2016), pp. 1457-1486.
\bibitem[9]{Weg17} Wegener, Patrick Hurwitz action in Coxeter groups and elliptic Weyl groups, Ph. D. Thesis, Universität Bielefeld (Germany) (2017) \url{https://pub.uni-bielefeld.de/record/2913106}.
\end{thebibliography}
\end{document}
